% ---------------------------------------------------------------------
% 02-main-result.tex
% The main theorem and its proof, plus the supporting lemmas the proof
% needs. State the theorem, prove the supporting results, then prove the
% theorem itself.
% ---------------------------------------------------------------------

\section{Main result}
\label{sec:main-result}

\begin{theorem}[Main theorem]
\label{thm:main}
Let $X$ be a locally compact Hausdorff space and let $\mu$ be a Radon
measure on $X$. Suppose $f \colon X \to \R$ is admissible. Then
\begin{equation}
\label{eq:main}
  \int_X f \dd\mu \;\le\; \mu(X)^{1/2}\left(\int_X f^2 \dd\mu\right)^{1/2}.
\end{equation}
Moreover, equality holds in \cref{eq:main} if and only if $f$ is
$\mu$-almost everywhere constant on the set where it is nonzero.
\end{theorem}

Before proving \cref{thm:main} we record the auxiliary estimate that
drives the argument.

\begin{proposition}[Auxiliary estimate]
\label{prop:aux}
Under the hypotheses of \cref{thm:main}, for every $\varepsilon > 0$
there exists a compactly supported admissible $g$ with
\begin{equation}
\label{eq:aux}
  \norm{f - g}_{L^2(\mu)} < \varepsilon .
\end{equation}
\end{proposition}

\begin{proof}
Approximate $f$ from below by simple functions, then truncate the
support to a set of finite measure. The truncation is admissible by
\cref{lem:uniform-bound}, and \cref{eq:aux} follows by monotone
convergence. This is the only place where $\sigma$-finiteness of
$\mu$ is used.
\end{proof}

% ---------------------------------------------------------------------
% figures/example-figure.tex
% A TikZ figure kept in its own file so the section sources stay
% readable. Include it from a section with % ---------------------------------------------------------------------
% figures/example-figure.tex
% A TikZ figure kept in its own file so the section sources stay
% readable. Include it from a section with % ---------------------------------------------------------------------
% figures/example-figure.tex
% A TikZ figure kept in its own file so the section sources stay
% readable. Include it from a section with \input{figures/example-figure}.
% ---------------------------------------------------------------------
\begin{figure}[htbp]
  \centering
  \begin{tikzpicture}[
      node distance=14mm,
      box/.style={draw=linkblue, thick, rounded corners=1.5pt,
                  inner sep=4pt, minimum height=8mm, align=center,
                  font=\small},
      arr/.style={-{Stealth[length=2mm]}, thick, draw=linkblue}
    ]
    \node[box] (a) {statement};
    \node[box, right=of a] (b) {estimate};
    \node[box, right=of b] (c) {conclusion};
    \draw[arr] (a) -- (b);
    \draw[arr] (b) -- (c);
  \end{tikzpicture}
  \caption{A sample diagram. Captions are complete sentences ending in
    a period, and the figure is referenced from the text as
    \cref{fig:sample}.}
  \label{fig:sample}
\end{figure}
.
% ---------------------------------------------------------------------
\begin{figure}[htbp]
  \centering
  \begin{tikzpicture}[
      node distance=14mm,
      box/.style={draw=linkblue, thick, rounded corners=1.5pt,
                  inner sep=4pt, minimum height=8mm, align=center,
                  font=\small},
      arr/.style={-{Stealth[length=2mm]}, thick, draw=linkblue}
    ]
    \node[box] (a) {statement};
    \node[box, right=of a] (b) {estimate};
    \node[box, right=of b] (c) {conclusion};
    \draw[arr] (a) -- (b);
    \draw[arr] (b) -- (c);
  \end{tikzpicture}
  \caption{A sample diagram. Captions are complete sentences ending in
    a period, and the figure is referenced from the text as
    \cref{fig:sample}.}
  \label{fig:sample}
\end{figure}
.
% ---------------------------------------------------------------------
\begin{figure}[htbp]
  \centering
  \begin{tikzpicture}[
      node distance=14mm,
      box/.style={draw=linkblue, thick, rounded corners=1.5pt,
                  inner sep=4pt, minimum height=8mm, align=center,
                  font=\small},
      arr/.style={-{Stealth[length=2mm]}, thick, draw=linkblue}
    ]
    \node[box] (a) {statement};
    \node[box, right=of a] (b) {estimate};
    \node[box, right=of b] (c) {conclusion};
    \draw[arr] (a) -- (b);
    \draw[arr] (b) -- (c);
  \end{tikzpicture}
  \caption{A sample diagram. Captions are complete sentences ending in
    a period, and the figure is referenced from the text as
    \cref{fig:sample}.}
  \label{fig:sample}
\end{figure}


\begin{proof}[Proof of \cref{thm:main}]
Apply Cauchy--Schwarz to the pairing $\ip{f}{\mathbf{1}_X}$:
\[
  \int_X f \dd\mu
  = \ip{f}{\mathbf{1}_X}
  \le \norm{f}_{L^2}\,\norm{\mathbf{1}_X}_{L^2}
  = \left(\int_X f^2 \dd\mu\right)^{1/2} \mu(X)^{1/2},
\]
which is \cref{eq:main}. Approximate $f$ by the compactly supported
$g$ produced by \cref{prop:aux} and pass to the limit to reduce the
equality statement to the standard equality case of Cauchy--Schwarz.
\end{proof}

\begin{corollary}
\label{cor:consequence}
If in addition $\mu(X) = 1$ then $\int_X f \dd\mu \le \norm{f}_{L^2}$.
\end{corollary}

\begin{proof}
Specialise \cref{thm:main} to $\mu(X) = 1$.
\end{proof}
